\section{Basics}

\begin{prop}[Bijection]
    The map \( f  \) is a bijection if and only if there exists \( g : B \to A  \) such that \( f \circ g   \) is the identity map on \( B  \) and \( g \circ f  \) is the identity map on \( A  \).
\end{prop}

\begin{definition}[Binary Relations, Equivalence Relations, Equivalence Classes, Partitions]
    \begin{enumerate}
        \item[(1)] A \textbf{binary relation} on a set \( A  \) is a subset \( R  \) of \( A \times A  \) and we write \( a \sim b  \) if \( (a,b)  \in R \).
        \item[(2)] The relation \( \sim  \) on \( A  \) is said to be:
            \begin{enumerate} 
                \item[(a)] \textbf{Reflexive} if \( a \sim a  \) for all \( a \in A  \);
                \item[(b)] \textbf{Symmetric} if \( a \sim b  \), then \( b \sim a  \) for all \( a,b \in A \);
                \item[(c)] \textbf{Transitive} if \( a \sim b  \) and \( b \sim c  \) implies \( a \sim c  \) for all \( a,b,c \in A  \).
            \end{enumerate}
    \end{enumerate}
\end{definition}

\begin{prop}[Isomorphism Properties]
    If \( \varphi : G \to H  \) is an isomorphism, then 
    \begin{enumerate}
        \item[(a)] \( | G |  = | H |  \)
        \item[(b)] \( G  \) is abelian if and only if \( H  \) is abelian.
        \item[(c)] for all \( x \in G  \), \( | x  |  = | \varphi(x) |  \).
    \end{enumerate}
\end{prop}


